\chapter{Conclusion and Future Work}
\label{ch:conclusion}

% ===========================================================================
\section{Summary of Work}
\label{sec:summary}

This thesis presented the design and implementation of a single-lane \gls{vpu}
integrated with a five-stage pipelined RISC-V scalar core, targeting \gls{asic} tapeout
under the \texttt{rv32im\_zicsr\_zve32x\_zvl128b} instruction set architecture with
VLEN=128 bits.

The hardware hierarchy was designed from scratch in synthesisable SystemVerilog. The
scalar core implements the RV32IM base \gls{isa} with a five-stage pipeline (IF, ID,
EX, MEM, WB), EX$\rightarrow$MEM$\rightarrow$WB forwarding, load-use hazard stalling,
and a dedicated dispatch path for vector instructions through a four-entry instruction
queue. The VPU is a separate clock-synchronous subsystem with its own nine-state
\gls{fsm} (\texttt{vproc\_fsm.sv}) that drives a single execution lane containing seven
functional sub-units: an adder, a multiplier, a shifter, a logic unit, a comparator,
a min/max unit, and a reduction unit. A 48-bit structured control bus produced by the
vector decoder (\texttt{vproc\_vdecoder.sv}) passes through the instruction queue and
drives all downstream units without any re-decoding, providing a clean single-decode
pipeline architecture.

The VPU supports \gls{sew} widths of 8, 16, and 32 bits, \gls{lmul} values of 1, 2, 4,
and 8, and the full arithmetic, logical, shift, compare, reduction, and vector
load/store instruction set of the Zve32x profile. Masked operations are supported with
per-element granularity; compare instructions produce packed mask results atomically
through the \texttt{ST\_FINAL\_MASKING} \gls{fsm} state.

A synthesisable UART peripheral mapped at \texttt{0xFF000000} enables real-data
end-to-end testing. Bare-metal firmware in RISC-V assembly implements the ITU-R BT.601
RGB-to-grayscale conversion algorithm using \texttt{vmulhu.vx} at SEW=8, exploiting
the upper-byte extraction semantics to implement the fixed-point multiply-and-divide
in a single vector instruction per colour channel. The firmware was compiled using the
\texttt{riscv64-unknown-elf-gcc} toolchain under MSYS2 and loaded into the IMEM model
via the \texttt{elf2hex\_lines.py} converter.

All modules were verified in ModelSim. The instruction regression testbench
(\texttt{tb\_vproc\_all\_instr.sv}) achieved 172/172 directed tests passing. System
integration was verified by the \texttt{tb\_uart\_echo.sv} testbench (16/16 pixel
comparisons passing, ACK \texttt{0xAA} confirmed) and by \texttt{tb\_uart\_lena\_assert.sv}
(16384/16384 Lena image pixels matching the BT.601 reference, all eleven assertion
categories passing). The design was synthesised on the Intel Cyclone~V DE10-Standard
board, consuming 10,438 \glspl{alm} and achieving an $F_\text{max}$ of 47.18~MHz in
the fast-corner timing model.

% ===========================================================================
\section{Contributions}
\label{sec:contributions_ch7}

The following six original contributions are claimed by this work:

\begin{enumerate}[label=\arabic*.]

  \item \textbf{Complete Zve32x-compliant RTL implementation in SystemVerilog.}
    A full implementation of the RVV Zve32x instruction subset was delivered in
    synthesisable SystemVerilog. All modules comply with the ASIC design rules:
    synchronous active-low reset throughout, no \texttt{initial} blocks or
    \texttt{\$display} in synthesisable paths, no combinational latches, and all
    widths parameterised as functions of the VLEN, SEW, and LMUL parameters.
    This work is, to the best of the author's knowledge, the first publicly documented
    single-lane Zve32x RTL implementation that explicitly documents tapeout-viability
    constraints and verification results.

  \item \textbf{Pipelined scalar core with VPU dispatch stall and FIFO-based decoupling.}
    A five-stage RV32IM pipeline was designed with two independent hazard mechanisms:
    a data-hazard forwarding network for scalar RAW dependencies, and an occupancy-based
    stall that halts the scalar IF/ID stages when the VPU instruction queue is full. The
    four-entry synchronous \gls{fifo} decouples the scalar and vector pipelines, allowing
    scalar instructions to advance while the VPU executes multi-cycle vector operations.

  \item \textbf{Per-element serial reduction with busy-based FSM handshake.}
    The \texttt{vproc\_reduction} module was redesigned from a combinational adder tree
    to a single-adder serial accumulator. Each element is added to the running accumulator
    on successive clock cycles, eliminating the $O(\log_2 N)$-deep carry chain. This
    redesign improved the design's $F_\text{max}$ from 16.44~MHz to 47.18~MHz (a
    $2.87\times$ improvement) and reduced the ALM count from 13,998 to 10,438 (a 25\%
    reduction). The \texttt{reduction\_busy} signal provides a synchronous handshake to
    the FSM so that the result-write cycle is correctly timed without additional dead cycles.

  \item \textbf{Comprehensive 172-test directed verification suite.}
    The \texttt{tb\_vproc\_all\_instr.sv} testbench provides 172 directed test cases
    covering every opcode category in the Zve32x subset: integer ALU in all three
    operand forms (VV, VX, VI), reverse subtract, logic, shifts, min/max with signed
    and unsigned comparisons, packed-mask compare, all four VMULH variants in both VV
    and VX forms, widening operations (six variants), LMUL=2 multi-register group
    operations, and all five reduction operations. This suite provides a regression
    gate that can be re-run after any \gls{rtl} modification to confirm that no
    regression was introduced.

  \item \textbf{End-to-end UART testbench with pixel-accurate BT.601 verification.}
    The \texttt{tb\_uart\_echo.sv} and \texttt{tb\_uart\_lena\_assert.sv} testbenches
    exercise the full hardware stack at pin level, driving a modelled UART receiver and
    observing the UART transmitter. The system is verified against a reference BT.601
    implementation in simulation; 100\% accuracy (zero error on all 16384 pixels of
    the Lena test image) is achieved, demonstrating that the fixed-point coefficient
    approximation is exact for all integer pixel values in [0, 255].

  \item \textbf{FPGA prototype on the DE10-Standard achieving 47.18~MHz at 10,438 ALMs.}
    The full design was synthesised, placed, and routed on the Intel Cyclone~V
    5CSXFC6D6F31C6 device. The achieved $F_\text{max}$ of 47.18~MHz, while below
    the 50~MHz target in the slow-corner model, represents a substantial improvement
    over the initial 16.44~MHz after the serial reduction redesign. The prototype
    validates that the design is fully synthesisable with modern FPGA tooling and
    provides a concrete area reference point (10,438 ALMs) for \gls{asic} cell-count
    estimation.

\end{enumerate}

% ===========================================================================
\section{Limitations}
\label{sec:limitations_ch7}

Several limitations remain in the current implementation:

\begin{description}

  \item[Slide instructions not implemented.]
    The \texttt{vslide1up.vx} and \texttt{vslide1down.vx} instructions, which shift
    a vector register's elements up or down by one position and insert a scalar value,
    are part of the Zve32x specification but have not been decoded or executed.
    These instructions are required for several important algorithms, including FIR
    filters using the overlap-save technique and certain recurrence computations.

  \item[Behavioral SRAM model.]
    The data memory is modelled as a \texttt{logic} array inside \texttt{dmem\_sync.sv}.
    While this is functionally correct in simulation and synthesises to FPGA block RAM,
    it is not a tapeout-ready SRAM macro. An ASIC implementation requires a standard-cell
    SRAM macro with defined access time, power ports, and boundary conditions.

  \item[Slow-corner timing violation.]
    The design fails to meet the 50~MHz target in the Cyclone~V slow-corner model by
    4.436~ns. The critical path is the 32-bit carry chain in the scalar ALU.
    While the design is functionally correct (the fast-corner model passes), an ASIC
    tapeout at 50~MHz would require this path to be closed under the worst-case
    process, voltage, and temperature corner.

  \item[No formal verification.]
    All verification was performed using directed simulation in ModelSim. No bounded
    model checking, property checking, or formal equivalence checking between the RTL
    and any reference model has been performed. Formal methods would provide stronger
    correctness guarantees, particularly for the decode logic and the CSR update path.

  \item[Asynchronous reset in lsu.sv.]
    The scalar load/store unit (\texttt{rtl/riscv/lsu.sv}) was identified as containing
    at least one register with an asynchronous reset path. This is inconsistent with
    the synchronous active-low reset policy applied to all other modules and must be
    corrected before tapeout.

  \item[Single-lane execution only.]
    The design achieves a peak throughput of one element per clock cycle. Multi-lane
    execution (two or four lanes operating in parallel) would proportionally increase
    throughput at the cost of additional area, but the current architecture was
    intentionally constrained to a single lane to minimise complexity during the
    initial implementation.

\end{description}

% ===========================================================================
\section{Future Work}
\label{sec:future}

The current design provides a solid, verified foundation for further development
along several axes.

\subsection{ASIC Synthesis and Physical Design}

The most direct continuation of this work is completing the \gls{asic} tapeout flow:

\begin{itemize}
  \item \textbf{Logic synthesis with Yosys or Synopsys Design Compiler.}
    Running the RTL through an open-source synthesis tool (Yosys with the
    \texttt{synth\_xilinx} or \texttt{synth} target) would provide a gate-level
    netlist and a technology-mapped area estimate that is more accurate than the
    FPGA ALM count. A commercial tool such as Synopsys Design Compiler with the
    target PDK would provide the production-quality netlist.
  \item \textbf{Static timing analysis at target frequency.}
    After synthesis, PrimeTime or OpenSTA should be run with a complete set of
    timing constraints covering all clock domains, I/O delays, and multicycle paths.
    The scalar ALU carry chain identified in this work would be the first path to
    address.
  \item \textbf{SRAM macro replacement.}
    The behavioral DMEM model should be replaced by a standard-cell SRAM macro
    (such as those generated by OpenRAM for the SkyWater SKY130 PDK or
    GlobalFoundries GF180MCU PDK). The dual-port access pattern of \texttt{dmem\_sync}
    requires a two-port or pseudo-dual-port SRAM with one read and one write port.
  \item \textbf{Physical design (floorplan, P\&R, DRC, LVS).}
    With a synthesis-clean netlist and SRAM macros, the design can be taken through
    OpenROAD or Cadence Innovus for floorplanning, placement, clock tree synthesis,
    routing, design-rule checking, and layout-versus-schematic verification.
    A final GDSII file suitable for multi-project wafer submission would be the
    deliverable.
\end{itemize}

\subsection{Formal and Property-Based Verification}

The simulation-based approach used in this thesis provides high confidence but cannot
prove the absence of bugs in unexplored corner cases. Future work should include:

\begin{itemize}
  \item \textbf{Bounded model checking} of the VPU \gls{fsm} to prove that no
        invalid state is reachable from reset and that all \texttt{busy}/\texttt{done}
        handshakes terminate within a bounded number of cycles.
  \item \textbf{Formal equivalence checking} between the RTL and the gate-level
        netlist to confirm that synthesis did not introduce any logical transformation.
  \item \textbf{Property-based verification} of the decoder, proving that every
        valid Zve32x encoding produces a unique, well-formed control word and that
        no two distinct encodings produce the same control word.
\end{itemize}

\subsection{Instruction Set Extensions}

The following instructions are within the Zve32x or closely related profiles and
would broaden the application coverage of the processor:

\begin{itemize}
  \item \textbf{Slide instructions} (\texttt{vslide1up.vx}, \texttt{vslide1down.vx},
        \texttt{vslideup.vx}, \texttt{vslidedown.vx}): required for prefix-sum,
        FIR, and IIR filter kernels. The implementation requires a barrel-shift network
        in the VRF read path.
  \item \textbf{Gather/scatter} (\texttt{vrgather.vv}, \texttt{vrgatherei16.vv}):
        enables permutation-based algorithms and look-up-table acceleration.
  \item \textbf{Segment load/store} (\texttt{vlseg}, \texttt{vsseg}): needed for
        interleaved (AoS) memory layouts such as packed RGB image data.
  \item \textbf{Single-width floating-point} (RVV with Zve32f extension): adding
        a 32-bit IEEE 754 floating-point datapath to the execution lane would support
        inference workloads that use \texttt{float32} activations.
\end{itemize}

\subsection{Multi-Lane Scaling}

The single-lane architecture of this design processes one element per cycle. Extending
to $L$ parallel lanes would process $L$ elements per cycle, proportionally increasing
throughput at the cost of $L\times$ the datapath area. A clean multi-lane extension
would require:

\begin{itemize}
  \item Partitioning the \gls{vrf} across $L$ banks so that each lane has a dedicated
        read/write port. With VLEN=128 and SEW=32, a single register group contains
        four elements, so four lanes is a natural choice that keeps each lane's bank
        width at 32 bits.
  \item Distributing the element-to-lane mapping consistently across all sub-units.
  \item Extending the \gls{vlsu} to generate $L$ DMEM requests per cycle.
  \item Updating the reduction unit to perform an $L$-input partial reduction across
        lanes before feeding the serial accumulator.
\end{itemize}

An implementation using two lanes would double throughput for most arithmetic
instructions and achieve approximately 4.2 cycles per pixel for the BT.601 kernel at
SEW=8.

\subsection{Out-of-Order VPU Execution with Scoreboard}

The current design stalls the scalar pipeline whenever the VPU instruction queue is
full or when a \texttt{vsetvli} result is pending. An out-of-order capability would
allow the scalar core to continue dispatching subsequent scalar instructions past a
pending vector stall, using a scoreboard to track which scalar registers are in use
as implicit VPU operands. This is particularly beneficial for scalar code that follows
a vector store instruction, which currently stalls for 16 or more cycles at SEW=8.

\subsection{Clock Gating and Power Optimisation}

The VRF data-path registers are the dominant source of dynamic power in the design.
Inserting clock-gating cells controlled by the element-enable signals (\texttt{vrf\_we})
would suppress switching in unused lanes during masked operations and during scalar
instructions when the VPU is idle. The tapeout checklist in \texttt{CLAUDE.md}
identifies this as a required item before final submission.

% ===========================================================================
\section{Closing Remarks}
\label{sec:closing}

This project set out to produce a standards-compliant RISC-V vector processor that is
not merely functionally correct in simulation but is also designed to the disciplined
engineering standard required for eventual \gls{asic} tapeout. The completed design
demonstrates that both goals can be achieved within the constraints of a single-person
academic project: full Zve32x instruction coverage, 172/172 regression passing,
pixel-accurate image processing verification, and a working FPGA prototype.

The central technical lesson is that timing closure and area efficiency for a vector
processor depend critically on the implementation of the reduction unit. A na\"ive
combinational tree appears straightforward but produces a combinational path whose
depth grows linearly with the vector length, making it incompatible with any practical
clock frequency target. The serial accumulator pattern — processing one element per
cycle under \gls{fsm} control — not only solves the timing problem but also reduces
area, because the synthesiser instantiates a single adder rather than a tree of
adders. This observation is transferable to any SIMD or vector design that includes
a reduction operation.

The open-source RVV ecosystem has matured considerably over the past three years.
Freely available toolchains, reference simulators (Spike), and the detailed RVV 1.0
specification provide an unusually solid foundation for academic hardware development.
The \texttt{rv32im\_zicsr\_zve32x\_zvl128b} profile is a well-scoped minimum viable
vector ISA that can be implemented and verified in a thesis-scale project while still
being directly useful for embedded machine-learning inference and image processing.

The work documented in this thesis provides a complete, well-verified starting point
for a future multi-project wafer submission. With the SRAM model replacement, the
scalar ALU timing fix, and the missing slide instructions completed, the design would
be ready for synthesis under an open-source PDK and subsequent physical design closure.
That tapeout remains the long-term goal, and this thesis is one rigorous step along
that path.
